\documentclass[10pt,a4paper]{amsart}
\usepackage[margin=19mm]{geometry}
\usepackage{amsmath,amssymb,mathtools}
\usepackage[scr=rsfs,cal=euler]{mathalfa}
\usepackage{xfrac}
\usepackage[unicode,hidelinks,bookmarks=false]{hyperref}
\newcommand{\ie}{i.e.\ }
\newcommand{\cf}{cf.\ }
\newcommand{\resp}{respectively\ }
\newcommand{\Subsection}{\subsection}
\providecommand{\nobreakdash}{\nobreak\hbox{-}}
\setlength{\emergencystretch}{2em}
\multlinegap0pt
\usepackage{bm}
\usepackage{bbm}
\usepackage{dsfont}
\usepackage{xspace}
\usepackage{cancel}
\usepackage{mathtools}
\usepackage{amsthm}
\makeatletter
\@ifundefined{theo}{\newtheorem{theo}{Theo}}{}
\@ifundefined{lem}{\newtheorem{lem}[theo]{Lemma}}{}
\makeatother
\newcommand{\E}{{\mathbb E}}
\newcommand{\N}{{\mathbb N}}
\newcommand{\R}{{\mathbb R}}
\newcommand{\dom}{{\mathrm{dom}}\;}
\newcommand{\ind}{\mathbf{1}}
\renewcommand{\P}{{\mathbb P}}
\title[Csisz\'ar indices and interpolating copulas]{Csisz\'ar indices and interpolating copulas}
\author{Cristina Butucea, Jean-Fran\c{c}ois Delmas, Anne Dutfoy, Antoine Schoonaert}
\date{Source: arXiv:2603.29884v1 (CC BY 4.0)}
\begin{document}
\maketitle

\noindent\emph{Source and licence.} Csisz\'ar indices and interpolating copulas. Version arXiv:2603.29884v1 (CC BY 4.0), distributed under the Creative Commons Attribution 4.0 International licence, as recorded in the arXiv licence metadata for this exact version and shown on the abstract page https://arxiv.org/abs/2603.29884v1. Full rights notice: licenses/arxiv-2603.29884-CC-BY-4.0.txt.\par\smallskip

\noindent\textit{Editorial scope.} Counted unit: the Lem whose source label is \texttt{lem:equality\_jensen} in the article (source0.tex lines 1678--1693, source label \texttt{lem:equality\_jensen}), delivered together with the article's own complete original proof of it (source0.tex lines 1707--1758, one \texttt{proof} environment of 52 lines). No mathematics is written, repaired or re-derived: statement and proof are the article's own text line for line. The proof is detached from the statement in the source: the article states the result at lines 1678--1693 and prints its proof at lines 1707--1758, and the proof environment's own header names the statement, so the association is the article's own. Required context retained uncounted: none. Every definition, display and label of those lines is retained unchanged. Omitted from this package and not redistributed: 1--1677, 1694--1706, 1759--2393. Nothing inside a retained excerpt is omitted or altered: every retained body is delivered exactly as the article's lines are written, so no omission receipt of any kind accompanies this package. Citations: the retained excerpts carry no citation command at all, so no bibliography entry of the article is transcribed and no reference is redistributed. Reference adaptation: none was needed and none was made; every reference inside the delivered unit resolves inside it, so no unresolved reference or citation is produced. Printed slips kept literal: no printed word, symbol, hypothesis or number of the counted statement or of its proof was altered. Layout and class: the article's own class is \texttt{article} with packages including geometry, changepage, hyperref, amsmath, amsfonts, amssymb, dsfont, bbm, stmaryrd, amsthm, graphicx, subcaption, tikz, pgfplots; the wrapper is the lane's pinned \texttt{amsart} editorial wrapper at 10pt on A4, which restates the theorem-like environments the retained text needs and adds the article's own macro definitions that the retained text needs (\texttt{\textbackslash E}, \texttt{\textbackslash N}, \texttt{\textbackslash P}, \texttt{\textbackslash R}, \texttt{\textbackslash dom}, \texttt{\textbackslash ind}), with extra wrapper packages (bm, bbm, dsfont, xspace, cancel, mathtools). The delivered numbering is the unit's own. Assets: no retained span contains a figure environment, a graphics-inclusion command, a PDF-page inclusion, a table or a TikZ picture; no image file is copied into this package, the package declares no asset dependency and no image file is redistributed with it. Licence: version arXiv:2603.29884v1 of this article is distributed under the Creative Commons Attribution 4.0 International licence, as recorded in the arXiv licence metadata for this exact version and shown on the abstract page https://arxiv.org/abs/2603.29884v1. A full rights notice is delivered at licenses/arxiv-2603.29884-CC-BY-4.0.txt. Characters: every retained excerpt is pure ASCII, so the delivered engine, XeLaTeX with its default Latin Modern text font, reports no missing character.
\begin{lem}[Jensen's inequality]
\label{lem:equality_jensen}
Let X be an $\bar \R$-valued  random variable. Let $\varphi$ be a proper
convex function defined on $\R$ (and by continuity on $\bar \R$). Assume
that $X$  and $\varphi(X)$ are  quasi-integrable (this is  in particular
the case if $X$ is integrable). Then, we have:
\begin{equation}
  \label{eq:def-proper-cvx-E}
 \varphi(\E[X]) \leq \E[\varphi(X)].
\end{equation}


Furthermore if $X$ is integrable and $\varphi$ is strictly convex at
$\E[X]$, then the inequality~\eqref{eq:def-proper-cvx-E} is an equality
if and only if a.s.\ $X =\E[X]$.
\end{lem}

\begin{proof}[Proof of Lemma~\ref{lem:equality_jensen}]
%  The proof  is standard, but we  give a short proof  for the interested reader.
  First notice that to prove~\eqref{eq:def-proper-cvx-E}, it is
  enough to  consider the case  $\E[\varphi(X)]\in [-\infty , + \infty )$
  and   $X$   non  constant.


  We then claim  that $\varphi$ is continuous  at $\E[X]$.  Indeed,
  since this is  the case by convention if $\E[X]=\pm\infty  $, we shall
  only consider the case $X$ integrable  and thus a.s.\ $X\in \R$. Then,
  since a.s.\ $\varphi(X)<+\infty $, we deduce that a.s.\ $X$ belongs to
  $\dom(\varphi)$. The latter being convex, $X$ non constant and $\E[X]$
  finite,  we   deduce  that   $\E[X]$  belongs   to  the   interior  of
  $\dom(\varphi)$. As the function $\varphi$ is convex, it is continuous
  in  the interior  of  its  domain. In conclusion,  the  function $\varphi$  is
  continuous at $\E[X]\in \bar \R$.

 Let $(X_n, n\in \N^*)$ be independent random
  variables distributed as $X$. Set $\bar X_n=n^{-1} \sum_{k=1}^n X_k$
  for $n\in \N^*$.
  Since $X$ and $\varphi(X)$ are quasi-integrable, the
  strong law of large number gives that a.s.\ 
  $$
  \lim_{n\rightarrow \infty
  } \bar X_n= \E[X] \text{ and }\lim_{n\rightarrow \infty
  } n^{-1} \sum_{k=1}^n \varphi(X_k)= \E[\varphi(X)].
  $$
Since $\varphi$ is convex, we also get that $\varphi(\bar X_n) \leq
n^{-1} \sum_{k=1}^n \varphi(X_k)$. Since $\varphi$ is continuous at
$\E[X]$,
we deduce that a.s.\ $\lim_{n\rightarrow \infty
} \varphi(\bar X_n) = \varphi(\E[X])$;
and thus~\eqref{eq:def-proper-cvx-E}
holds.


Assume  that $X$  is  integrable  and $\varphi$  is  strictly convex  at
$\E[X]\in \R$.  This implies that
$\E[X]\in \dom (\varphi)$  and as~\eqref{eq:def-proper-cvx-E} is an
equality, we deduce that a.s.\ $X \in \dom  (\varphi)$.  We shall
assume  that $X$  is  not  constant. Since  $X$  is  integrable and  non
constant,  this  implies  that  $\E[X]$   belongs  to  the  interior  of
$\dom ( \varphi)$.   Thus,  there  exists an  affine  function  $g$  which
coincides with $\varphi$  at $\E[X]$ and such that $\varphi  \geq g$ and
$\varphi(t)>g(t)$ for all $t>\E[X]$ or  for all $t<\E[X]$.  Without loss
of  generality, we  can  assume that  $\varphi>g$  on $A=\{t>  \E[X]\}$.
Since   $X$  is   not  constant,   we  get   $\P(X\in  A)>0$   and  thus
$\E[\varphi(X)]=        \E[\varphi(X)\ind_{\{X\in         A\}}]        +
\E[\varphi(X)\ind_{\{X\not \in A\}}]> \E[g(X)]=g(\E[X])=\varphi(\E[X])$.
So if~\eqref{eq:def-proper-cvx-E}  is an equality, we  deduce that a.s.\
$X=\E[X]$. \hfill \qed
\end{proof}

\end{document}
